\makeatletter
\def\input@path{{./Graphics/}} % Définition du chemin d'entrée pour les graphiques
\makeatother

\documentclass[12pt,a4paper]{article} % Classe de document

% Chargement des packages essentiels
\usepackage[utf8]{inputenc} % Encodage UTF-8
\usepackage[french]{babel} % Support de la langue française
\usepackage{graphicx, amsthm, mathtools, amsfonts, amssymb, physics, longtable, mathrsfs} % Packages pour les graphiques et mathématiques

% Divers Box en couleur avec le package tcolorbox
\usepackage[many]{tcolorbox}
\usepackage{varioref}
%
% Définition avec caption et label=df:
\newtcbtheorem[number within=section]{definition}{Définition}
{fonttitle=\bfseries\upshape, fontupper=\slshape, arc=0mm, colback=blue!3,colframe=blue!75!black,separator sign dash,breakable,beamer}{df}
%
\newtcbtheorem[number within=section]{example}{Exemple}
{fonttitle=\bfseries\upshape,fontupper=\slshape, arc=0mm, colback=green!5!white,
colframe=green!55!black,separator sign dash,breakable,beamer}{exp}

\usepackage{float} % Gère le placement des objets flottants
\floatplacement{figure}{H} % Forces les figures à être placées exactement à l'endroit où elles apparaissent dans le texte

\usepackage[svgnames,x11names,rgb,table]{xcolor} % Définit les couleurs par noms
\usepackage{array, multicol, wrapfig, lmodern, multirow, booktabs, chemformula} % Diverses améliorations de mise en forme
\usepackage[detect-all=true, group-minimum-digits=4]{siunitx} % Pour les unités SI

\usepackage[bf, textfont={footnotesize,it}]{caption} % Personnalisation des légendes
\usepackage{subcaption}

% Configuration des marges de la page
\usepackage[margin=2cm, a4paper]{geometry}

\usepackage{fancyvrb} % Remplacement du mode verbatim qui permet une meilleure mise en forme
%\usepackage[Export]{adjustbox} % Utilisé pour contraindre les images à une taille maximale
%\adjustboxset{max size={0.9\linewidth}{0.9\paperheight}} % Taille maximale des images

% Configuration pour hyperref
\usepackage{hyperref} 
\hypersetup{
 colorlinks=true, % Colorise les liens
 breaklinks=true, % Permet le retour à la ligne dans les liens trop longs
 allcolors=Blue3, % Couleur des liens
 pdfauthor={Nana Engo}, % Auteur du document PDF
 pdfsubject={Chimie quantique et VQE}, % Sujet du document
 pdftitle={PHY4268 - Qubit Hamiltonian V25.04}, % Titre du document
 bookmarksopen=true % Ouvre les signets PDF par défaut
}

\usepackage{cleveref} % Améliore la gestion des références croisées

\frenchbsetup{StandardLayout=true} % Configuration du layout français
\numberwithin{equation}{section} % Numérotation des équations par section
\numberwithin{figure}{section} % Numérotation des figures par section
\numberwithin{table}{section} % Numérotation des tables par section

\setcounter{secnumdepth}{3} % Pour inclure les sous-sections
\renewcommand{\thesection}{\arabic{section}} % Numérotation des sections
%
\newcommand{\tcb}[1]{\textcolor{blue}{#1}} % textcolor{blue}{#1}
\newcommand{\tcr}[1]{\textcolor{red}{#1}} % textcolor{red}{#1}
% Titre du document
\title{\vspace*{4cm}
\Huge\tcb{Chapter 2 - Hamiltonien Qubit}
\vspace*{5cm}}

\author{
\begin{tabular}{lr}
 \textbf{Nana Engo} & \textbf{Tchapet Njafa}\\
 {\small serge.nana-engo@facsciences-uy1.cm} &\hspace*{2em}
 {\small jean-pierre.tchapet@facsciences-uy1.cm}
\end{tabular}\\[3em] % Espace entre les auteurs et la filiation
\textit{Quantum Simulation Group}\\
 \textit{Department of Physics} \\ 
 \textit{Faculty of Science} \\ 
 \textit{University of Yaoundé I}
}
\date{\vspace*{3cm}April 2025} % Date du document

\begin{document}
\maketitle % Affiche le titre, les auteurs et la date

\newpage
\section{Introduction}

L'objectif de ce Chapitre est de présenter brièvement comment on établi une correspondance entre les opérateurs d'échelle fermioniques (Hamiltonien en seconde quantification ou Hanmiltonien fermionique) et les opérateurs de Pauli (Hamiltonien Qubit), à travers les transformations de Jordan-Wigner (JWT), de parité (PT) ou de Bravyi-Kitaev (BKT) (voir les parties marquées en rouge sur la Figure 1 présentant la VQE). De façon spécifique, à la fin de ce chapitre, l'apprenant doit
* avoir compris le principe de l'encodage d'un problème de structure électronique en seconde quantification sur un calculateur quantique;
* être capable 
	* de transformer l'Hamiltonien en 1ère quantification en Hamiltonien en seconde quantification (ou Hamiltonien fermionique);
    * transformer un Hamiltonien fermionique en Hamiltonien Qubit.

![VQE_Flowchart](./Graphics/VQE_Flowchart_Tuto2.jpeg)

\section{Hamiltonien moléculaire fermionique}

\subsection{Génération des intégrales moléculaires}

L'Hamiltonien  est exprimé dans la base des solutions de la méthode HF, également appelées Orbitales Moléculaires (OM) :
$$
\mathtt{H}_{elec}=\sum_{pq} h_{pq} a^{\dagger}_p a_q + 
\frac{1}{2} \sum_{pqrs} h_{pqrs}  a^{\dagger}_p a^{\dagger}_q a_r  a_s
$$
avec 
* les **intégrales à 1 électron**
$$
h_{pq} = \int \phi^*_p(r) \left( -\frac{1}{2} \nabla^2 - \sum_{I} \frac{Z_I}{|\mathbf{r}_i-\mathbf{R}_I|} \right)   \phi_q(r)d\mathbf{r}
$$
qui décrivent l’énergie cinétique des électrons individuels et leurs interactions Coulombiennes avec les champs électriques du noyau ;
* et **intégrales à 2 électrons**
$$
h_{pqrs} = \int \frac{\phi^*_p(r_1)  \phi^*_q(r_2) \phi_r(r_2)  \phi_s(r_1)}{|\mathbf{r}_1-\mathbf{r}_2|}d\mathbf{r}_1d\mathbf{r}_2,
$$
décrivent les interactions Coulombiennes entre les électrons;
* $a^\dagger_p$ l'**opérateur de création** d'un électron sur le spin-orbite $p$;
* $a_q$ l'**opérateur d'annihilation** d'un électron sur sur le spin-orbite $q$;
* $a^{\dagger}_p a_q $ est l'**opérateur d'excitation**, qui excite un électron du spin-orbite occupé $\phi_q$ vers le spin-orbite inoccupé $\phi_p$.

Les MO ($\phi_u$) peuvent être occupés ou virtuels (inoccupés). Un MO peut contenir 2 électrons. Ces MO sont les solutions de la méthode HF, où chaque l´electron est soumis au champ moyen créé par les autres électrons. Cependant, dans ce qui suit, nous travaillons en fait avec des orbitales de spin qui sont associées à un spin up ($\alpha$) d'électron spin down ($\beta$). Ces deux spins sont également communément désignés par $\alpha$ et $\beta$, respectivement. Ainsi, les orbitales de spin peuvent contenir un électron ou être inoccupées.

Dans le tableau ci-dessous resume les correspondances première-seconde quantification des divers éléments du Hamiltonien moléculaire. 
$$
\begin{array}{|l|l|l|}\hline
\textbf{Approximation BO}  &  \textbf{Combinaison en 2e quantification} &  \textbf{Description}
\\\hline
-\frac12\sum_i\nabla^2_i & -\frac12\sum_{p,q}\langle p|\nabla^2_p|q\rangle a^\dagger_p a_q & \mathtt{T}_e \\\hline
 -\sum_{i,I}\frac{Z_I}{|\mathbf{r}_i-\mathbf{R}_I|} &
 -\sum_{p,q} \langle p|\frac{Z_I}{|\mathbf{r}_p-\mathbf{R}_I|}|q\rangle a^\dagger_p a_q
 & \mathtt{U}_{en}\\\hline
\sum_{i,j>i}\frac{1}{|\mathbf{r}_i-\mathbf{r}_j|} & 
\sum_{p,q,r,s}\langle pq \big|\frac{1}{|\mathbf{r}_i-\mathbf{r}_j|}\big|rs\rangle a^\dagger_p a^\dagger_q a_r a_s & 
 \mathtt{U}_{ee}\\\hline
\end{array}
$$

*Un des avantages apporté par le formalisme de seconde quantification est que la propriété d'anti-symétrie des fonctions d'état pour l'échange de fermions identiques, prise en compte manuellement dans une approche de première quantification en considérant les déterminants de Slater, est automatiquement imposée par le relations d'anti-commutation des opérateurs de création et d'annihilation. Cependant, le prix à payer est de travailler avec un nombre non fixe de particules, qui reste de toute façon une quantité conservée.*

### <a id='toc1_3_2_'></a>[Répresentation nombre d'occupation](#toc0_)

En seconde quantification, le déterminant de Slater est représenté par un **vecteur nombre d'occupation**
\begin{equation*}
|f\rangle =|f_{M-1},\dots,f_{i-1},f_i,f_{i+1},\dots, f_0\rangle, 
\end{equation*}
avec
- $f_i=1$ lorsque le spin-orbital $\phi_i$ est occupé et donc présent dans le déterminant de Slater;
- $f_i = 0$ lorsque le $\phi_i$ spin-orbital est vide et donc absent du déterminant de Slater.

- Le vecteur $|f_i\rangle$ est appelé **vecteur nombre d'occupation de l'orbite fermionique i**, parce que
\begin{align*}
& |0\rangle=\begin{pmatrix}1\\0\end{pmatrix}:\text{état vide}, & |1\rangle=\begin{pmatrix}0\\1\end{pmatrix}:\text{état occupé}.
\end{align*}

Dans la deuxième représentation de quantification, au lieu de poser la question *Quel électron est dans quel état?*, nous posons la question *combien de particules y a-t-il dans chaque état?*.

### <a id='toc1_3_3_'></a>[Opérateurs fermioniques](#toc0_)

Les opérateurs fermioniques $a^\dagger_i$ et $a_i$ obéissent aux relations d'anticommutation fermioniques suivantes:
\begin{align*}
&\{a_p,a^\dagger_q\} = a_pa^\dagger_q + a^\dagger_q a_p = \delta_{pq}, 
&\{a_p, a_q \} = \{a^\dagger_p, a^\dagger_q\} = 0 .
\end{align*}
La nature fermionique des électrons implique que les fonctions d'état à plusieurs électrons doivent être antisymétriques par rapport à l'échange de particules. Cela se reflète dans la manière dont les opérateurs de création fermionique et d'annhilation agissent sur les déterminants $|f\rangle$:
\begin{equation*}
\begin{aligned}
&a_i|f_{M-1},\dots,f_{i-1},f_i,f_{i+1},\dots, f_0\rangle = \delta_{f_i,1}
(-1)^{\sum_{m<i}f_m}|f_{M-1},\dots,f_{i-1},f_i\oplus 1,f_{i+1},\dots, f_0\rangle,\\
&a^\dagger_i|f_{M-1},\dots,f_{i-1},f_i,f_{i+1},\dots, f_0\rangle = \delta_{f_i,0}
(-1)^{\sum_{m<i}f_m}|f_{M-1},\dots,f_{i-1},f_i\oplus 1,f_{i+1},\dots, f_0\rangle`,  \\
&a|0\rangle=a^\dagger|1\rangle=0.
\end{aligned}
\end{equation*}
- Le terme de phase $p_i=(-1)^{\sum_{m<i}f_m}$ désigne la parité qu'impose l'anti-symétrie d'échange des fermions.
    - $p_i=-1$ si le nombre d'électrons est impair dans ces orbitales de spin,
    - $p_i= 1$ si le nombre d'électrons est pair dans ces orbitales de spin.

- Le symbole $\oplus$ représente l'addition modulo 2, i.e.,  $1 \oplus 1=0 $ et $0 \oplus 1=1$.

- L'opérateur de création $a^\dagger_i$ met un électron dans une orbitale inoccupée *i* ou dans $|f_i\rangle$.

- L'opérateur d'annihilation $a_i$ supprime un électron dans une orbitale occupée *i* ou en $|f_i\rangle$.

- $a_i^2 =(a^\dagger_i)^2= 0$ pour tout $i$. On ne peut pas créer ou anéantir un fermion dans le même mode deux fois.

Par example,
\begin{align*} 
&a^\dagger_1 |0\rangle_1 = |1\rangle_1,&&a^\dagger_1 |1\rangle_1 =(a^\dagger_i)^2|0\rangle_1= 0,&& a_1 |1\rangle_1 =|0\rangle_1,&a_1 |0\rangle_1 =a_1^2|1\rangle_1=0. 
\end{align*}
Notez qu'ici $a^\dagger_1|1\rangle_1=0$ et $a_1 |0\rangle_1=0$ signifie le vecteur zéro et non $|0\rangle_1$. 

En utilisant également de tels opérateurs, nous pouvons exprimer
\begin{equation*}
|0\rangle |1\rangle |1\rangle |0\rangle = a^\dagger_1a^\dagger_2 |0\rangle^{\otimes 4}.
\end{equation*}

 L'opérateur d'occupation orbitale est donné par
\begin{align*}
&n_i = a^\dagger_i a_i, &n_i |f_{M-1},\dots,f_i,\dots, f_0\rangle= f_i |f_{M-1},\dots,f_i,\dots, f_0\rangle,
\end{align*}
et il compte le nombre d'électrons dans une orbitale donnée, c'est-à-dire
\begin{align*} 
&n_i |0\rangle_i = 0, &n_i |1\rangle_i = |1\rangle_i. 
\end{align*}

### <a id='toc1_3_4_'></a>[Exercise](#toc0_)

In **Exemple d'illustration** above, we have considered a fictitious system described by spin orbitals $|A_\uparrow\rangle,\,|A_\downarrow\rangle,\,|B_\uparrow\rangle,\,|B_\downarrow\rangle$ labelled as
\begin{align*}
 &|A_\uparrow\rangle=|00\rangle=|\mathbf{0}\rangle,  &&|A_\downarrow\rangle=|01\rangle=|\mathbf{1}\rangle,
&&|B_\uparrow\rangle=|10\rangle=|\mathbf{2}\rangle,  &|B_\downarrow\rangle=|11\rangle=|\mathbf{3}\rangle.
\end{align*}
We assume that the Hartree-Fock state for this system **has both electrons occupying the $|A\rangle$ orbitals**. We store the occupations of the spin orbitals $|A_\uparrow\rangle,\,|A_\downarrow\rangle,\,|B_\uparrow\rangle,\,|B_\downarrow\rangle$ which we order as 
$|f_{B_\downarrow},f_{B_\uparrow},f_{A_\downarrow},f_{A_\uparrow}\rangle$, with $f_i=0,1$ (in the fermionic Fock space).

1. Write down, in the fermionic Fock space, the Hartree-Fock state $|\Psi_{\rm HF}\rangle$, that is, the corresponding of the antisymmetrized Slater deter-
minant obtained in **Exemple d'illustration**.
$$
|\Psi_{\rm HF}\rangle = |0011\rangle .
$$
1. Write down, in the fermionic Fock space, the $s_z = 0$ state function.
$$
|\Psi_{\rm HF}\rangle = \alpha|0011\rangle + \beta|1100\rangle + |1001\rangle + |0101\rangle .
$$


\section{Hamiltonien Qubit}

\subsection{Principe de l'optimisation géométrique}

\begin{definition}{Optimisation de la structure géométrique}{}
L'\textbf{optimisation de la structure géométrique}, souvent appelée simplement optimisation géométrique, optimisation moléculaire ou relaxation, est une étape primordiale en modélisation moléculaire. Elle désigne le processus par lequel la géométrie d'une molécule est ajustée de manière itérative, dans le but de trouver une structure correspondant à une énergie électronique minimale.
\end{definition}

\subsubsection{Pourquoi optimiser la géométrie ?}

La \textbf{recherche de la géométrie moléculaire d'énergie minimale} est essentielle car, en vertu des principes fondamentaux de la nature, les molécules tendent à adopter les conformations les plus stables, c'est-à-dire celles qui minimisent leur énergie totale. Ces structures optimisées correspondent généralement à des configurations stables : 
\begin{itemize}
\item soit un \textbf{minimum local} sur la surface d'énergie potentielle (représentant une espèce moléculaire stable), 
\item soit un \textbf{état de transition}, représentant une configuration de haute énergie entre deux minima, crucial pour comprendre les mécanismes réactionnels.
\end{itemize}

La \autoref{fig:pes3d} illustre de manière visuelle ce concept de minimisation de l'énergie. Imaginez la surface d'énergie potentielle comme un paysage montagneux complexe. Les molécules stables correspondent aux vallées de ce paysage (minima d'énergie), tandis que les chemins réactionnels entre ces vallées passent par des cols de montagne (états de transition). L'optimisation géométrique est alors analogue à la recherche des points les plus bas dans ce paysage, ou à l'identification des cols qui permettent de passer d'une vallée à une autre.

\begin{figure}[h!]
\centering
\includegraphics[width=0.35\textwidth]{PES_TS.png} 
\caption{\textbf{Surface d'Énergie Potentielle 3D et optimisation géométrique.} On peut visualiser une Surface d'Énergie Potentielle (PES) à deux dimensions comme un paysage 3D. Les zones bleues foncées représentent les régions de basse énergie (vallées), correspondant aux \textbf{minima} de la PES et aux structures moléculaires stables. Les zones plus claires représentent des régions de plus haute énergie. Le chemin jaune tracé sur la surface illustre un \textbf{chemin réactionnel} typique, reliant deux minima en passant par un \textbf{état de transition (TS)}, situé au point col de la PES. L'optimisation géométrique consiste à partir d'une géométrie initiale et à se déplacer itérativement sur la PES, \textit{en descendant la pente} du gradient, jusqu'à atteindre un minimum local d'énergie.}
\label{fig:pes3d}
\end{figure}


\subsubsection{Surface d'Énergie Potentielle (PES) et optimisation}

Comme nous l'avons introduit précédemment, le concept central pour l'optimisation géométrique est la \textbf{Surface d'Énergie Potentielle (PES)}. Dans l'approximation de Born-Oppenheimer, nous considérons les noyaux comme des particules ponctuelles fixes. L'énergie totale de la molécule, $E(\mathbf{x})$, devient alors une fonction des seules coordonnées nucléaires $\mathbf{x}$ (positions des noyaux). Cette fonction $E(\mathbf{x})$ définit la PES.

L'optimisation géométrique revient mathématiquement à résoudre le problème stationnaire $\nabla E(\mathbf{x}) = 0$, c'est-à-dire à trouver les points où le gradient de l'énergie (la pente de la PES) s'annule. Ces points stationnaires correspondent aux minima (structures stables) et aux points cols (états de transition) sur la PES. Les coordonnées nucléaires $\mathbf{x}$ obtenues après optimisation définissent alors la \textbf{géométrie d'équilibre} de la molécule, ou la géométrie de l'état de transition.

\subsection{Ingrédients pour l'optimisation géométrique}

En pratique, la réalisation d'une optimisation géométrique nécessite de définir et de calculer plusieurs éléments clés :
\begin{enumerate}
 \item \textbf{Coordonnées moléculaires initiales.} Point de départ de l'optimisation. Elles peuvent être estimées à partir de données expérimentales (cristallographie, spectroscopie), de l'intuition chimique, de banques de données moléculaires, ou générées par des outils de construction moléculaire (logiciels de dessin moléculaire, outils comme RDKit ou Openbabel).

 \item \textbf{Choix du système de coordonnées.} Comme discuté dans la section précédente, le choix entre coordonnées cartésiennes et coordonnées internes influence l'efficacité de l'optimisation. Les coordonnées internes sont souvent privilégiées pour leur pertinence chimique et leur meilleur découplage des mouvements.

 \item \textbf{Calcul de l'énergie $E(\mathbf{x})$ et du gradient $\mathbf{g}(\mathbf{x}) = \nabla E(\mathbf{x})$.} Pour chaque géométrie moléculaire testée lors de l'optimisation, il est nécessaire de calculer l'énergie électronique $E(\mathbf{x})$ et son gradient $\mathbf{g}(\mathbf{x})$. Ces calculs sont réalisés à l'aide d'une méthode de chimie quantique (Hartree-Fock, DFT, méthodes post-Hartree-Fock, méthodes semi-empiriques ou Tight-Binding). Le gradient $\mathbf{g}(\mathbf{x})$ représente la force exercée sur chaque atome et indique la direction de la plus forte variation de l'énergie.

 \item \textbf{Hessien (optionnel, mais souvent utilisé pour les optimisations avancées).} La matrice Hessienne, $\mathbf{H}(\mathbf{x})$, est la matrice des dérivées secondes de l'énergie par rapport aux coordonnées nucléaires. Elle décrit la courbure de la PES au point $\mathbf{x}$. Le Hessien est utilisé dans les algorithmes d'optimisation plus avancés (comme la méthode de Newton-Raphson) pour accélérer la convergence et pour caractériser la nature des points stationnaires (minima, états de transition).
 
  \textbf{Importance du Hessien pour la caractérisation des points stationnaires.} Afin de vérifier si un point stationnaire, c'est-à-dire un point sur la surface d'énergie potentielle où le gradient s'annule, correspond à un minimum local (structure stable) ou à un état de transition (point col), il est nécessaire de calculer les dérivées secondes de l'énergie par rapport aux déplacements nucléaires. Cette matrice de toutes les dérivées secondes possibles est appelée \textbf{matrice Hessienne}.

  Une fois calculée, la matrice Hessienne peut être diagonalisée pour obtenir les \textbf{constantes de force harmoniques}, les \textbf{fréquences vibrationnelles} et les \textbf{modes normaux de vibration} de la molécule au voisinage du point stationnaire. L'analyse des valeurs propres (fréquences vibrationnelles) du Hessien permet de déterminer la nature du point stationnaire :
  \begin{itemize}
    \item \textbf{Minimum local.} Toutes les valeurs propres du Hessien sont positives (fréquences vibrationnelles réelles).
    \item \textbf{État de transition.} Une (et une seule) valeur propre du Hessien est négative (fréquence vibrationnelle imaginaire), correspondant au mouvement le long de la coordonnée de réaction.
    \item \textbf{Maximum d'ordre supérieur.} Plus d'une valeur propre négative (rare en chimie moléculaire, généralement instable).
  \end{itemize}

 \item \textbf{Algorithme d'optimisation.} Procédure algorithmique pour mettre à jour itérativement les coordonnées atomiques et se déplacer sur la PES vers des régions de plus basse énergie. Les algorithmes courants incluent la descente de gradient, la méthode de Newton-Raphson et ses variantes (BFGS, etc.).

 \item \textbf{Critères de convergence.} Conditions définissant quand l'optimisation est considérée comme terminée. Les critères de convergence typiques portent sur :
 \begin{itemize}
  \item Le gradient maximal (force maximale sur un atome) qui doit être inférieur à un seuil.
  \item Le déplacement maximal des atomes entre deux itérations qui doit être inférieur à un seuil.
  \item La variation de l'énergie entre deux itérations qui doit être inférieure à un seuil.
  \item La convergence de l'énergie elle-même, cad l'énergie doit se stabiliser et ne plus varier significativement.
 \end{itemize}
\end{enumerate}

\subsection{Fonctionnement d'une optimisation géométrique}

La \autoref{fig:geooptworkflow} illustre les étapes typiques d'une optimisation de la géométrie moléculaire.

\begin{figure}[h!]
\centering
\includegraphics[width=0.6\textwidth]{Geo_opt.png} 
\caption{\textbf{Workflow d'une optimisation géométrique moléculaire.} Ce diagramme illustre le processus itératif typique d'une optimisation géométrique. L'optimisation débute avec une géométrie moléculaire initiale et une estimation du Hessien. À chaque itération, un calcul SCF (Self-Consistent Field) est réalisé pour obtenir l'énergie et le gradient. Un test de convergence est effectué sur le gradient et le déplacement des atomes. Si la convergence n'est pas atteinte, le Hessien et la géométrie sont mis à jour, et le processus est répété. Une fois la convergence atteinte, la géométrie optimisée est sauvegardée.}
\label{fig:geooptworkflow}
\end{figure}

Le processus d'optimisation se déroule de manière itérative :

\begin{enumerate}
 \item \textbf{Initialisation (étape 1 de la \autoref{fig:geooptworkflow}).} L'optimisation commence avec une \textbf{géométrie moléculaire initiale} fournie par l'utilisateur et une estimation du \textbf{Hessien} initial (souvent une matrice Hessienne approximée ou un Hessien calculé pour une méthode simple).
 
 \item \textbf{Calcul SCF et Gradient (étape 2 de la \autoref{fig:geooptworkflow}).} À chaque itération, un calcul de structure électronique (typiquement un calcul SCF dans le cas de la méthode Hartree-Fock ou DFT) est réalisé pour la géométrie moléculaire courante. Ce calcul fournit l'\textbf{énergie électronique} $E$ et le \textbf{gradient de l'énergie} $\nabla E$ (les forces sur les atomes) pour cette géométrie.
 
 \item \textbf{Test de Convergence (étape 3 de la \autoref{fig:geooptworkflow}).} Les critères de convergence (sur le gradient, le déplacement atomique, la variation d'énergie) sont évalués. Si tous les critères sont satisfaits, l'optimisation est considérée comme \textbf{convergée}, et la géométrie courante est considérée comme la \textbf{géométrie optimisée} (correspondant à un minimum local ou un état de transition).
 
 \item \textbf{Mise à jour du Hessien (étape 4 de la \autoref{fig:geooptworkflow}) et de la Géométrie (étape 5 de la \autoref{fig:geooptworkflow}).} Si la convergence n'est pas atteinte, le \textbf{Hessien} est mis à jour (par exemple, en utilisant des méthodes de mise à jour du Hessien comme BFGS) pour mieux approximer la courbure de la PES au voisinage du minimum. Ensuite, un \textbf{algorithme d'optimisation} (par exemple, descente de gradient, BFGS) utilise le gradient et (éventuellement) le Hessien pour prédire une \textbf{nouvelle géométrie moléculaire} avec une énergie plus basse. Les coordonnées atomiques sont mises à jour en conséquence.
 
 \item \textbf{Itération suivante.} Le processus retourne à l'étape 2 avec la nouvelle géométrie moléculaire, et le cycle itératif se poursuit jusqu'à convergence.

\end{enumerate}

\subsubsection{Minima locaux vs. Minimum global}

Il est important de souligner que les algorithmes d'optimisation géométrique convergent généralement vers le \textbf{minimum local} le plus proche de la géométrie initiale. Bien que l'objectif idéal soit de trouver le \textbf{minimum global} (la structure la plus stable de la molécule), il n'existe pas de méthode algorithmique universelle pour garantir la localisation du minimum global pour des molécules complexes. En pratique, plusieurs stratégies peuvent être utilisées pour explorer différentes régions de la PES et augmenter les chances de trouver le minimum global, telles que :

\begin{itemize}
 \item Démarrer l'optimisation à partir de différentes géométries initiales, générées par différentes méthodes (construction manuelle, dynamique moléculaire, algorithmes de recherche conformationnelle).
 \item Utiliser des algorithmes d'optimisation plus sophistiqués, capables d'échapper aux minima locaux (par exemple, "basin hopping", algorithmes génétiques).
 \item Combiner l'optimisation géométrique avec des méthodes de recherche conformationnelle plus larges (par exemple, dynamique moléculaire simulée, algorithmes génétiques).
\end{itemize}

\subsection{Avantages de l'optimisation orbitale}

Le terme \textbf{optimisation orbitale}" est parfois utilisé de manière interchangeable avec l'optimisation géométrique, car dans les méthodes de structure électronique comme Hartree-Fock et DFT, l'énergie moléculaire dépend à la fois des positions nucléaires (géométrie) et des orbitales électroniques. L'optimisation géométrique implique donc une optimisation conjointe de la géométrie nucléaire et, de manière implicite, des orbitales électroniques, qui s'adaptent à chaque nouvelle géométrie nucléaire via le calcul SCF.

L'optimisation orbitale dans les méthodes de structure électronique offre les avantages suivants :
\begin{enumerate}
 \item \textbf{Énergies variationnelles.} Les énergies obtenues sont variationnelles en ce qui concerne les rotations orbitales. Cela signifie qu'il n'y a pas besoin de considérer la réponse orbitale lors du calcul des gradients nucléaires, ce qui simplifie le calcul des forces.

 \item \textbf{Calcul simplifié des propriétés.} Les propriétés moléculaires peuvent être calculées plus facilement, car il n'y a aucune contribution de réponse orbitale aux matrices de densité.

 \item \textbf{Suppression des pôles parasites.} Les pôles parasites qui peuvent apparaître dans les fonctions de réponse pour des méthodes approximées comme la théorie des clusters couplés (CC) peuvent être supprimés grâce à l'optimisation orbitale.

 \item \textbf{Meilleure description des ruptures de symétrie.} L'optimisation orbitale peut améliorer la description des problèmes de rupture de symétrie qui peuvent se produire dans certaines méthodes de structure électronique.
\end{enumerate}

\subsection{Conclusion sur l'optimisation géométrique}

L'optimisation géométrique est une étape fondamentale et incontournable dans la modélisation moléculaire. Elle permet de déterminer les structures moléculaires stables, correspondant aux minima de la surface d'énergie potentielle (PES). Le choix de l'algorithme d'optimisation, des critères de convergence, et du niveau de théorie quantique sont des aspects importants à considérer pour obtenir des résultats précis et fiables. Comprendre le principe de l'optimisation géométrique et les algorithmes associés est essentiel pour tout chercheur souhaitant utiliser la chimie quantique computationnelle pour étudier les propriétés et la réactivité des molécules.

\section{Conclusion}

Dans ce chapitre, nous avons posé les bases essentielles pour la représentation et la manipulation des molécules dans le cadre de la chimie quantique computationnelle. Nous avons exploré les différentes manières de décrire la géométrie moléculaire, en mettant en évidence les avantages et les inconvénients des \textbf{coordonnées cartésiennes}, des \textbf{coordonnées internes (Matrice Z)}, et du format \textbf{SMILES}. Vous devriez maintenant être en mesure de choisir le système de coordonnées le plus adapté à une application donnée, et de comprendre les principes fondamentaux des formats de fichiers utilisés en chimie moléculaire numérique.

Nous avons également introduit en détail la notion d'\textbf{optimisation géométrique}, une étape essentielle pour déterminer les structures moléculaires stables. Vous avez découvert le concept de \textbf{Surface d'Énergie Potentielle (PES)}, le rôle du \textbf{gradient} et du \textbf{Hessien}, ainsi que l'\textbf{algorithme itératif} général d'optimisation et les \textbf{critères de convergence} associés. Vous comprenez maintenant pourquoi l'optimisation géométrique est indispensable pour obtenir des résultats pertinents et comparables aux données expérimentales.

En conclusion, ce chapitre vous a fourni les outils conceptuels et pratiques pour :
\begin{itemize}
 \item \textbf{Représenter une molécule} de manière appropriée pour les calculs de chimie quantique, en choisissant le système de coordonnées et le format de fichier adaptés.
 
 \item \textbf{Comprendre et mettre en œuvre une optimisation géométrique} pour trouver les structures moléculaires d'énergie minimale, en interprétant les résultats en termes de Surface d'Énergie Potentielle.
\end{itemize}

Ces compétences sont fondamentales pour la suite de votre apprentissage en chimie quantique computationnelle. Dans les chapitres suivants, nous nous appuierons sur ces bases pour explorer des méthodes de calcul plus avancées et pour étudier les propriétés moléculaires avec une précision accrue.
\end{document}
